\ifdefined\gkver\else\def\gkver{student}\fi
\documentclass[\gkver]{gaokaozhenti}
\begin{document}

\chapter{1992年全国}

\begin{enumerate}
\item
%% number: 1
%% paperName: 1992年全国高考第 1 题 2 分
%% typeId: 1
%% type: 单选题
%% chapter: 电场
%% point: 电场线
%% method: 无
%% score: 2
%% degree: 850
%% duplicateId: 0
%% body:
如图所示，$Q$ 是带正电的点电荷，$P_{1}$ 和 $P_{2}$ 为其电场中的两点。若 $E_{1}$、$E_{2}$ 为 $P_{1}$、$P_{2}$ 两点的电场强度的大小，$\varphi_{1}$、$\varphi_{2}$ 为 $P_{1}$、$P_{2}$ 两点的电势，则\xzanswer{A}
\onepicture[w=4.5cm]{figs/01.png}
\fourchoices[answer=A]
{$E_{1} > E_{2}$，$\varphi_{1} > \varphi_{2}$}
{$E_{1} > E_{2}$，$\varphi_{1} < \varphi_{2}$}
{$E_{1} < E_{2}$，$\varphi_{1} > \varphi_{2}$}
{$E_{1} < E_{2}$，$\varphi_{1} < \varphi_{2}$}

%% answer: A
%% memo:
\memoanswer{
因 $Q$ 是带正电的点电荷，由电场强度的表达式为 $E=\frac{kQ}{r^{2}}$，所以 $E_{1}>E_{2}$；又沿电场线电势降低的最快，所以 $\varphi_{1}>\varphi_{2}$，故 A 正确。
}

\item
%% number: 2
%% paperName: 1992年全国高考第 2 题 2 分
%% typeId: 1
%% type: 单选题
%% chapter: 气体性质
%% point: 盖吕萨克定律
%% method: 无
%% score: 2
%% degree: 850
%% duplicateId: 0
%% body:
一定质量的理想气体，在压强不变的条件下，体积增大。则\xzanswer{A}
\fourchoices[answer=A]
{气体分子的平均动能增大}
{气体分子的平均动能减少}
{气体分子的平均动能不变}
{条件不够，无法判定气体分子平均动能的变化}

%% answer: A
%% memo:
\memoanswer{
由等压变化得 $\frac{V_{1}}{T_{1}}=\frac{V_{2}}{T_{2}}$，因体积增大，所以温度升高，因此气体分子的平均动能增大，故 A 正确。
}

\item
%% number: 3
%% paperName: 1992年全国高考第 3 题 2 分
%% typeId: 1
%% type: 单选题
%% chapter: 机械波
%% point: 波的图象
%% method: 无
%% score: 2
%% degree: 750
%% duplicateId: 0
%% body:
a、b 是一条水平的绳上相距为 $l$ 的两点。一列简谐横波沿绳传播，其波长等于 $\frac{2}{3}l$，当 a 点经过平衡位置向上运动时，b 点\xzanswer{C}
\fourchoices[answer=C]
{经过平衡位置向上运动}
{处于平衡位置上方位移最大处}
{经过平衡位置向下运动}
{处于平衡位置下方位移最大处}

%% answer: C
%% memo:
\memoanswer{
由题意可知 $\lambda = \frac{2}{3}l$，即 $l = \frac{3}{2}\lambda$，a、b 两点相距 1.5 个波长。如果波传播方向为从 a 向 b，则如图所示，b 经过平衡位置向下运动，故 C 正确；

\onepicture[w=6cm]{figs/03_an1.png}

如果波传播方向为从 b 向 a，则如图所示，b 经过平衡位置向下运动，故 C 仍正确。

\onepicture[w=6cm]{figs/03_an2.png}
}

\item
%% number: 4
%% paperName: 1992年全国高考第 4 题 2 分
%% typeId: 1
%% type: 单选题
%% chapter: 曲线运动
%% point: 人造地球卫星
%% method: 无
%% score: 2
%% degree: 850
%% duplicateId: 0
%% body:
两颗人造地球卫星，都在圆形轨道上运行，它们的质量相等，轨道半径之比 $\frac{r_{1}}{r_{2}} = 2$，则它们动能之比 $\frac{E_{1}}{E_{2}}$ 等于\xzanswer{C}
\fourchoices[answer=C]
{2}
{$\sqrt{2}$}
{$\frac{1}{2}$}
{4}

%% answer: C
%% memo:
\memoanswer{
由万有引力提供向心力得 $G\frac{Mm}{r^{2}}=m\frac{v^{2}}{r}$，得动能 $E_{k}=G\frac{Mm}{2r}$，所以 $\frac{E_{k1}}{E_{k2}}=\frac{r_{2}}{r_{1}}=\frac{1}{2}$，故 C 正确。
}

\item
%% number: 5
%% paperName: 1992年全国高考第 5 题 2 分
%% typeId: 1
%% type: 单选题
%% chapter: 原子和原子核
%% point: 原子核式结构
%% method: 无
%% score: 2
%% degree: 850
%% duplicateId: 0
%% body:
卢瑟福 $\alpha$ 粒子散射实验的结果\xzanswer{C}
\fourchoices[answer=C]
{证明了质子的存在}
{证明了原子核是由质子和中子组成的}
{说明原子的全部正电荷和几乎全部质量都集中在一个很小的核上}
{说明原子中的电子只能在某些不连续的轨道上运动}

%% answer: C
%% memo:
\memoanswer{
$\alpha$ 粒子散射实验的结果是：绝大多数 $\alpha$ 粒子穿过金箔后直进，少数 $\alpha$ 粒子偏转，极少数 $\alpha$ 粒子偏转角度较大或被弹回，说明 $\alpha$ 粒子遇到了原子中电量和质量非常集中且体积很小的部分，该实验说明原子的全部正电荷和几乎全部质量都集中在一个很小的核上，故 C 正确；卢瑟福就是通过对该结果的深入分析和研究，提出了原子的核式结构学说。
}

\item
%% number: 6
%% paperName: 1992年全国高考第 6 题 2 分
%% typeId: 1
%% type: 单选题
%% chapter: 运动和力
%% point: 牛顿第二定律
%% method: 无
%% score: 2
%% degree: 750
%% duplicateId: 0
%% body:
如图，位于水平地面上的质量为 $M$ 的小木块，在大小为 $F$、方向与水平方向成 $\alpha$ 角的拉力作用下沿地面作加速运动。若木块与地面之间的滑动摩擦系数为 $\mu$，则木块的加速度为\xzanswer{D}
\onepicture[w=5cm]{figs/06.png}
\fourchoices[answer=D]
{$\frac{F}{M}$}
{$\frac{F\cos \alpha }{M}$}
{$\frac{F\cos \alpha  - \mu Mg}{M}$}
{$\frac{F\cos \alpha  - \mu (Mg - F\sin \alpha )}{M}$}

%% answer: D
%% memo:
\memoanswer{
以小木块为研究对象，它受重力 $Mg$、地面的摩擦力 $f$、地面的支持力 $N$、外力 $F$ 的作用。将 $F$ 沿水平和竖直方向分解，水平方向分力 $F_{1}=F\cos\alpha$，竖直方向分力 $F_{2}=F\sin\alpha$，小木块在竖直方向受力平衡，有 $F_{2}+N=Mg$，小木块在水平方向加速运动，有 $F_{1}-f=Ma$，又由摩擦定律得 $f=\mu N$，联立解得 $a=\frac{F\cos\alpha-\mu(Mg-F\sin\alpha)}{M}$，故 D 正确。
}

\item
%% number: 7
%% paperName: 1992年全国高考第 7 题 2 分
%% typeId: 1
%% type: 单选题
%% chapter: 电场
%% point: 带电粒子的偏转
%% method: 无
%% score: 2
%% degree: 550
%% duplicateId: 0
%% body:
如图，电子在电势差为 $U_{1}$ 的加速电场中由静止开始运动，然后射入电势差为 $U_{2}$ 的两块平行极板间的电场中，入射方向跟极板平行。整个装置处在真空中，重力可忽略。在满足电子能射出平行板区的条件下，下述四种情况中，一定能使电子的偏转角 $\theta$ 变大的是\xzanswer{B}
\onepicture[w=6cm]{figs/07.png}
\fourchoices[answer=B]
{$U_{1}$ 变大、$U_{2}$ 变大}
{$U_{1}$ 变小、$U_{2}$ 变大}
{$U_{1}$ 变大、$U_{2}$ 变小}
{$U_{1}$ 变小、$U_{2}$ 变小}

%% answer: B
%% memo:
\memoanswer{
电子在加速电场中运动时，由动能定理得 $qU_{1}=\frac{1}{2}mv^{2}$；在偏转电场中运动时，水平方向做匀速运动，有 $L=vt$，竖直方向做加速运动，有 $v_{y}=at=\frac{qU_{2}}{dm}t$，则 $\tan\theta=\frac{v_{y}}{v}=\frac{at}{v}=\frac{qU_{2}}{dm}\cdot\frac{L}{v^{2}}$，解得 $\tan\theta=\frac{U_{2}L}{2dU_{1}}$，其中 $L$、$d$ 分别表示平行极板的长度和板间距离。从上式可知：一定能使电子的偏转角 $\theta$ 变大的是 $U_{1}$ 变小、$U_{2}$ 变大，故 B 正确。
}

\item
%% number: 8
%% paperName: 1992年全国高考第 8 题 2 分
%% typeId: 1
%% type: 单选题
%% chapter: 电路
%% point: 动态分析
%% method: 无
%% score: 2
%% degree: 750
%% duplicateId: 0
%% body:
如图的电路中，电池的电动势为 $\varepsilon$，内阻为 $r$，$R_{1}$ 和 $R_{2}$ 是两个阻值固定的电阻。当可变电阻 $R$ 的滑片向 a 点移动时，通过 $R_{1}$ 的电流 $I_{1}$ 和通过 $R_{2}$ 的电流 $I_{2}$ 将发生如下的变化\xzanswer{C}
\onepicture[w=5cm]{figs/08.png}
\fourchoices[answer=C]
{$I_{1}$ 变大，$I_{2}$ 变小}
{$I_{1}$ 变大，$I_{2}$ 变大}
{$I_{1}$ 变小，$I_{2}$ 变大}
{$I_{1}$ 变小，$I_{2}$ 变小}

%% answer: C
%% memo:
\memoanswer{
当可变电阻 $R$ 的滑片向 a 点移动时，电阻 $R$ 变小，干路电流 $I$ 变大，由 $U=E-Ir$ 知，路端电压 $U$ 变小，又通过 $R_{1}$ 的电流 $I_{1}=\frac{U}{R_{1}}$，所以 $I_{1}$ 变小，则通过 $R_{2}$ 的电流 $I_{2}=I-I_{1}$，所以 $I_{2}$ 变大，故 C 正确。
}

\item
%% number: 9
%% paperName: 1992年全国高考第 9 题 2 分
%% typeId: 1
%% type: 单选题
%% chapter: 交流电
%% point: 交流电
%% method: 无
%% score: 2
%% degree: 650
%% duplicateId: 0
%% body:
交流发电机在工作时的电动势为 $e=\varepsilon_{0}\sin\omega t$，若将其电枢的转速提高 1 倍，其他条件不变，则其电动势变为\xzanswer{D}
\fourchoices[answer=D]
{$\varepsilon_{0}\sin2\omega t$}
{$2\varepsilon_{0}\sin2\omega t$}
{$\varepsilon_{0}\sin\frac{\omega t}{2}$}
{$2\varepsilon_{0}\sin\frac{\omega t}{2}$}

%% answer: D
%% memo:
\memoanswer{
电动势瞬时表达式为 $e=E_{0}\sin\omega t=NBS\omega\sin\omega t$，其中 $\omega=2\pi n$，当转速 $n$ 增大为原来的 2 倍时，$\omega$ 也变为原来的 2 倍，最大值也变为原来的 2 倍，其表达式应为 $e=NBS\cdot2\omega\sin2\omega t$，即 $2E_{0}\sin2\omega t$，故 D 正确。
}

\item
%% number: 10
%% paperName: 1992年全国高考第 10 题 2 分
%% typeId: 1
%% type: 单选题
%% chapter: 电路
%% point: 电路中的图象
%% method: 无
%% score: 2
%% degree: 750
%% duplicateId: 0
%% body:
两电阻 $R_{1}$、$R_{2}$ 的电流 $I$ 和电压 $U$ 的关系图线如图所示，可知两电阻的大小之比 $R_{1}:R_{2}$ 等于\xzanswer{A}
\onepicture[w=4.5cm]{figs/10.png}
\fourchoices[answer=A]
{1:3}
{3:1}
{1:$\sqrt{3}$}
{$\sqrt{3}$:1}

%% answer: A
%% memo:
\memoanswer{
由题意和图象知 $R = \cot\alpha$，故 $\frac{R_{1}}{R_{2}} = \frac{\cot60^\circ}{\cot30^\circ} = \frac{1}{3}$，故 A 正确。
}

\item
%% number: 11
%% paperName: 1992年全国高考第 11 题 2 分
%% typeId: 1
%% type: 单选题
%% chapter: 力物体的平衡
%% point: 摩擦力
%% method: 无
%% score: 2
%% degree: 550
%% duplicateId: 0
%% body:
如图，一木块放在水平桌面上，在水平方向共受到三个力即 $F_{1}$、$F_{2}$ 和摩擦力作用，木块处于静止状态。其中 $F_{1}=10\UN$、$F_{2}=2\UN$。若撤去力 $F_{1}$，则木块在水平方向受到的合力为\xzanswer{D}
\onepicture[w=6cm]{\includesvg{figs/11}}
\fourchoices[answer=D]
{$10\UN$，方向向左}
{$6\UN$，方向向右}
{$2\UN$，方向向左}
{零}

%% answer: D
%% memo:
\memoanswer{
撤去 $F_{1}$ 前，因 $F_{1}>F_{2}$，物块有向右运动趋势，静摩擦力向左，大小为 $8\UN$；若撤去力 $F_{1}$，则木块只受 $F_{2}$ 作用，所以摩擦力为 $2\UN$，方向向右，在水平方向受到的合力为 0，故 D 正确。
}

\item
%% number: 12
%% paperName: 1992年全国高考第 12 题 2 分
%% typeId: 1
%% type: 单选题
%% chapter: 运动和力
%% point: 动量和能量
%% method: 无
%% score: 2
%% degree: 650
%% duplicateId: 0
%% body:
如图所示的装置中，木块 B 与水平桌面间的接触是光滑的，子弹 A 沿水平方向射入木块后留在木块内，将弹簧压缩到最短。现将子弹、木块和弹簧合在一起作为研究对象（系统），则此系统在从子弹开始射入木块到弹簧压缩至最短的整个过程中\xzanswer{B}
\onepicture[w=6cm]{\includesvg{figs/12}}
\fourchoices[answer=B]
{动量守恒、机械能守恒}
{动量不守恒、机械能不守恒}
{动量守恒、机械能不守恒}
{动量不守恒、机械能守恒}

%% answer: B
%% memo:
\memoanswer{
系统在水平方向上，在子弹射入木块到它陷入木块中并相对木块的速度变为零的过程中，由于历时很短，以至于在这段极小时间内木块对弹簧的压缩量极小，从而弹簧的弹力及墙对弹簧的挤压力可忽略不计，因此在子弹射入木块、弹簧几乎尚未发生形变的这一碰撞过程中，系统的动量守恒。弹簧在压缩过程中，墙对弹簧向右的挤压力始终存在，并且逐渐变大，因此对于子弹、木块和弹簧组成的系统，外力的冲量在水平方向不为零，即系统的动量不守恒。由于在子弹射入木块过程中，相对木块做减速运动，有摩擦阻力做功导致子弹和木块的内能增加，此过程为完全非弹性碰撞，系统的机械能减小，即系统的机械能不守恒，故 B 正确。
}

\item
%% number: 13
%% paperName: 1992年全国高考第 13 题 2 分
%% typeId: 1
%% type: 单选题
%% chapter: 直线运动
%% point: 相遇问题
%% method: 图象法
%% score: 2
%% degree: 550
%% duplicateId: 0
%% body:
两辆完全相同的汽车，沿水平直路一前一后匀速行驶，速度均为 $v_{0}$，若前车突然以恒定的加速度刹车，在它刚停住时，后车以前车刹车时的加速度开始刹车。已知前车在刹车过程中所行的距离为 $s$，若要保证两辆车在上述情况中不相撞，则两车在匀速行驶时保持的距离至少应为\xzanswer{B}
\fourchoices[answer=B]
{$s$}
{$2s$}
{$3s$}
{$4s$}

%% answer: B
%% memo:
\memoanswer{
前车在刹车过程中所行的距离为 $s$，因汽车做匀减速运动，有 $s = \frac{v_{0}}{2}t$；在这段时间内，后车匀速运动，行驶的距离为 $v_{0}t = 2s$，然后后车刹车，后车在刹车过程中所行的距离也为 $s$，到达前车所停的位置，所以两车在匀速行驶时保持的距离至少应为 $2s$，故 B 正确。

也可用图象法解答，两车的 $v_{0}-t$ 图象如图所示。前车在刹车过程中所行的距离为三角形 $Otv_{0}$ 的“面积”，大小为 $s$，则后车在刹车过程中行驶的距离为三角形 $t2tA$ 的“面积”，大小也为 $s$，在前车刹车过程中，后车所行的距离为四边形 $OtAv_{0}$ 的“面积”，大小为 $2s$，所以两车在匀速行驶时保持的距离至少应为 $2s$，故 B 正确。

\onepicture[w=5cm]{figs/13_an.png}
}

\item
%% number: 14
%% paperName: 1992年全国高考第 14 题 4 分
%% typeId: 2
%% type: 多选题
%% chapter: 电路
%% point: 电容
%% method: 无
%% score: 4
%% degree: 650
%% duplicateId: 0
%% body:
平行板电容器的电容\xzanswer{BC}
\fourchoices[answer=BC]
{跟两极板间的距离成正比}
{跟充满极板间的介质的介电常数成正比}
{跟两极板的正对面积成正比}
{跟加在两极板间的电压成正比}

%% answer: BC
%% memo:
\memoanswer{
平行板电容器的电容跟两极板间的距离成反比，跟充满极板间的介质的介电常数成正比，跟两极板的正对面积成正比，故 BC 正确。
}

\item
%% number: 15
%% paperName: 1992年全国高考第 15 题 4 分
%% typeId: 2
%% type: 多选题
%% chapter: 电路
%% point: 串并联电路
%% method: 无
%% score: 4
%% degree: 750
%% duplicateId: 0
%% body:
如图所示，a、b、c、d 是滑线变阻器的 4 个接线柱。现把此变阻器串联接入电路中，并要求滑片 P 向接线柱 c 移动时，电路中的电流减小。则接入电路的接线柱可能是\xzanswer{CD}
\onepicture[w=5cm]{figs/15.png}
\fourchoices[answer=CD]
{a 和 b}
{a 和 c}
{b 和 c}
{b 和 d}

%% answer: CD
%% memo:
\memoanswer{
滑片 P 向接线柱 c 移动时，电路中的电流减小，原因是串联接入电路中变阻器的阻值变大，因此接入电路的接线柱只能是 b 和 c 或 b 和 d，故 CD 正确。
}

\item
%% number: 16
%% paperName: 1992年全国高考第 16 题 4 分
%% typeId: 2
%% type: 多选题
%% chapter: 磁场
%% point: 带电粒子在复合场中的运动
%% method: 无
%% score: 4
%% degree: 650
%% duplicateId: 0
%% body:
在图中虚线所围的区域内，存在电场强度为 $E$ 的匀强电场和磁感应强度为 $B$ 的匀强磁场。已知从左方水平射入的电子，穿过这区域时未发生偏转。设重力可忽略不计，则在这区域中的 $E$ 和 $B$ 的方向可能是\xzanswer{ABC}
\onepicture[w=4cm]{\includesvg{figs/16}}
\fourchoices[answer=ABC]
{$E$ 和 $B$ 都沿水平方向，并与电子运动的方向相同}
{$E$ 和 $B$ 都沿水平方向，并与电子运动的方向相反}
{$E$ 竖直向上，$B$ 垂直纸面向外}
{$E$ 竖直向上，$B$ 垂直纸面向里}

%% answer: ABC
%% memo:
\memoanswer{
$E$ 和 $B$ 沿水平方向并与电子运动方向平行（相同或相反），则洛伦兹力 $f=0$，只受电场力作用且电场力与电子速度方向共线，必不偏转，故 AB 正确；若 $E$ 向上、$B$ 向外，则电子运动中所受洛伦兹力向上，电场力向下，不偏转是有可能的，故 C 正确；而 D 是电场力向下，洛伦兹力也向下，电子受力方向与初速度方向不平行，必然偏转，故 D 错误。
}

\item
%% number: 17
%% paperName: 1992年全国高考第 17 题 4 分
%% typeId: 2
%% type: 多选题
%% chapter: 光学
%% point: 光的折射
%% method: 无
%% score: 4
%% degree: 850
%% duplicateId: 0
%% body:
红光与紫光相比\xzanswer{BD}
\fourchoices[answer=BD]
{在真空中传播时，紫光的速度比较大}
{在玻璃中传播时，红光的速度比较大}
{玻璃对红光的折射率较紫光的大}
{从玻璃到空气的界面上，红光的临界角较紫光的大}

%% answer: BD
%% memo:
\memoanswer{
在真空中传播时，红光与紫光的速度相同，等于 $c$；在玻璃中传播时，由 $v=\frac{c}{n}$ 知，因 $n_{\text{红}}<n_{\text{紫}}$，所以 $v_{\text{红}}>v_{\text{紫}}$，红光的速度比较大；玻璃对红光的折射率较紫光的小，从玻璃到空气的界面上，因临界角 $C=\arcsin\frac{1}{n}$，所以红光的临界角较紫光的大，故 BD 正确。
}

\item
%% number: 18
%% paperName: 1992年全国高考第 18 题 4 分
%% typeId: 2
%% type: 多选题
%% chapter: 力物体的平衡
%% point: 摩擦力
%% method: 无
%% score: 4
%% degree: 550
%% duplicateId: 0
%% body:
如图所示，位于斜面上的物块 $M$ 在沿斜面向上的力 $F$ 作用下，处于静止状态。则斜面作用于物块的静摩擦力的\xzanswer{ABCD}
\onepicture[w=5cm]{\includesvg{figs/18}}
\fourchoices[answer=ABCD]
{方向可能沿斜面向上}
{方向可能沿斜面向下}
{大小可能等于零}
{大小可能等于 $F$}

%% answer: ABCD
%% memo:
\memoanswer{
物块 $M$ 处于静止状态，设斜面倾角为 $\alpha$，沿斜面向上为正方向，设静摩擦力为 $f$，由物体的平衡条件，列力的平衡方程

\begin{enumerate}
	\item
	当 $mg\sin\alpha>F$ 时，$f>0$，即沿斜面向上，故 A 正确；
	\item
	当 $mg\sin\alpha<F$ 时，$f<0$，即 $f$ 的方向沿斜面向下，故 B 正确；
	\item
	当 $mg\sin\alpha=F$ 时，$f=0$，物块不受静摩擦力的作用，故 C 正确；
	\item
	当 $F=\frac{1}{2}mg\sin\alpha$ 时，$f=\frac{1}{2}mg\sin\alpha$，即 $f=F$，故 D 正确。
\end{enumerate}
}

\item
%% number: 19
%% paperName: 1992年全国高考第 19 题 4 分
%% typeId: 2
%% type: 多选题
%% chapter: 曲线运动
%% point: 线速度角速度
%% method: 无
%% score: 4
%% degree: 650
%% duplicateId: 0
%% body:
图中所示为一皮带传动装置，右轮的半径为 $r$，a 是它边缘上的一点。左侧是一轮轴，大轮的半径为 4$r$，小轮的半径为 2$r$。b 点在小轮上，到小轮中心的距离为 $r$。c 点和 d 点分别位于小轮和大轮的边缘上。若在传动过程中，皮带不打滑。则\xzanswer{CD}
\onepicture[w=5.5cm]{\includesvg{figs/19}}
\fourchoices[answer=CD]
{a 点与 b 点的线速度大小相等}
{a 点与 b 点的角速度大小相等}
{a 点与 c 点的线速度大小相等}
{a 点与 d 点的向心加速度大小相等}

%% answer: CD
%% memo:
\memoanswer{
因在传动过程中，皮带不打滑，则 a 点与 c 点的线速度大小相等，故 C 正确；由 $a=\frac{v^{2}}{r}$ 得 $a_{a}=2a_{c}$，由 $a=\omega^{2}r$ 得 $a_{d}=2a_{c}$，所以 $a_{a}=a_{d}$，故 D 正确。
}

\item
%% number: 20
%% paperName: 1992年全国高考第 20 题 3 分
%% typeId: 3
%% type: 填空题
%% chapter: 原子和原子核
%% point: 天然放射性
%% method: 无
%% score: 3
%% degree: 850
%% duplicateId: 0
%% body:
在中子、质子、电子、正电子、$\alpha$ 粒子中选出一个适当的粒子，分别填在下列核反应式的横线上：

\ce{^{238}_{92}U -> ^{234}_{90}Th +} \tkanswer[1.5cm]{$\ce{^{4}_{2}He}$}；

\ce{^{234}_{90}Th -> ^{234}_{91}Pa +} \tkanswer[1.5cm]{$\ce{^{0}_{-1}e}$}；

\ce{^{9}_{4}Be + ^{4}_{2}He -> ^{12}_{6}C +} \tkanswer[1.5cm]{$\ce{^{1}_{0}n}$}。

%% answer: $\ce{^{4}_{2}He}$，$\ce{^{0}_{-1}e}$，$\ce{^{1}_{0}n}$
\jdanswer{
$\ce{^{4}_{2}He}$，$\ce{^{0}_{-1}e}$，$\ce{^{1}_{0}n}$
}

%% memo:
\memoanswer{
首先由核电荷数守恒算出粒子的核电荷数从而确定粒子名称，然后再依据质量数守恒填上粒子的质量数。
}

\item
%% number: 21
%% paperName: 1992年全国高考第 21 题 3 分
%% typeId: 3
%% type: 填空题
%% chapter: 光学
%% point: 光电效应
%% method: 无
%% score: 3
%% degree: 650
%% duplicateId: 0
%% body:
已知铯的极限频率为 $4.545\times10^{14}\UHz$，钠的为 $6.000\times10^{14}\UHz$，银的为 $1.153\times10^{15}\UHz$，铂的为 $1.529\times10^{15}\UHz$。当用波长为 $0.375\Uum$ 的光照射它们时，可发生光电效应的是 \tkanswer[1.5cm]{铯、钠}。

%% answer: 铯、钠
\jdanswer{
铯、钠
}

%% memo:
\memoanswer{
由 $c=\lambda f$ 得 $f=\frac{c}{\lambda}=\frac{3\times10^{8}}{0.375\times10^{-6}}\UHz=8\times10^{14}\UHz$，此值大于铯和钠发生光电效应的极限频率，故答案为铯和钠。
}

\item
%% number: 22
%% paperName: 1992年全国高考第 22 题 3 分
%% typeId: 3
%% type: 填空题
%% chapter: 曲线运动
%% point: 竖直平面圆周运动
%% method: 无
%% score: 3
%% degree: 750
%% duplicateId: 0
%% body:
图中圆弧轨道 AB 是在竖直平面内的 $\frac{1}{4}$ 圆周，在 B 点，轨道的切线是水平的，一质点自 A 点从静止开始下滑，不计滑块与轨道间的摩擦和空气阻力，则在质点刚要到达 B 点时的加速度大小为 \tkanswer[1.5cm]{$2g$}，刚滑过 B 点时的加速度大小为 \tkanswer[1.5cm]{$g$}。
\onepicture[align=right, w=3.5cm]{figs/22.png}

%% answer: $2g$；$g$
\jdanswer{
$2g$；$g$
}

%% memo:
\memoanswer{
滑块在由 A 到 B 的圆周运动中，只受重力和支持力作用，由于支持力不做功，则在运动过程中滑块机械能守恒，由机械能守恒得 $mgh = \frac{1}{2}mv^{2}$，解得 $v = \sqrt{2gR}$，滑块在刚要到达 B 点时，向心加速度即为物体的加速度，则 $a = a_{\text{向}} = \frac{v^{2}}{R} = \frac{2gR}{R} = 2g$，当滑块滑过 B 点后，由于只受重力，则物体加速度即为重力加速度 $g$，即质点刚要到达 B 点时的加速度大小为 $2g$，刚滑过 B 点时的加速度大小为 $g$。
}

\item
%% number: 23
%% paperName: 1992年全国高考第 23 题 3 分
%% typeId: 3
%% type: 填空题
%% chapter: 电路
%% point: 电表改装
%% method: 无
%% score: 3
%% degree: 650
%% duplicateId: 0
%% body:
一量程为 $0.6\UA$ 的电流表，其刻度盘如上图所示。今在此电流表的两端间并联一电阻，其阻值等于该电流表内阻的 $\frac{1}{2}$，使之成为一新的电流表，则图示的刻度盘上的每一小格表示 \tkanswer[1.5cm]{$0.06\UA$}。
\onepicture[align=right, w=5cm]{figs/23.png}

%% answer: $0.06\UA$
\jdanswer{
$0.06\UA$
}

%% memo:
\memoanswer{
原来刻度盘上每一小格表示 $0.02\UA$。当并联一个阻值等于电流表内阻 $\frac{1}{2}$ 的电阻后，通过这个电阻的电流是通过电流表电流的 2 倍，总电流是只有电流表一条通路时电流的 3 倍，因此此时电流表每一小格表示 $0.06\UA$。
}

\item
%% number: 24
%% paperName: 1992年全国高考第 24 题 3 分
%% typeId: 3
%% type: 填空题
%% chapter: 直线运动
%% point: 游标卡尺和螺旋测微仪
%% method: 无
%% score: 3
%% degree: 850
%% duplicateId: 0
%% body:
在测定金属丝的直径时，螺旋测微器的读数如图所示。可知该金属丝的直径 $d =$ \tkanswer[1.2cm]{0.900}$\times 10^{-3}\Um$。
\onepicture[align=right, w=2.5cm]{figs/24.png}

%% answer: $0.900\times10^{-3}\Um$
\jdanswer{
$0.900\times10^{-3}\Um$
}

%% memo:
\memoanswer{
固定刻度读出 $0.5\Umm$，可动部分是 50 等分，每一等分表示 $\frac{0.5\Umm}{50} = 0.01\Umm$，因此，可动部分应该读出 $40.0 \times 0.01\Umm = 0.400\Umm$，故读数为 $0.5\Umm + 40.0 \times 0.01\Umm = 0.900\Umm = 0.900 \times 10^{-3}\Um$。
}

\item
%% number: 25
%% paperName: 1992年全国高考第 25 题 3 分
%% typeId: 3
%% type: 填空题
%% chapter: 力物体的平衡
%% point: 平衡综合
%% method: 无
%% score: 3
%% degree: 650
%% duplicateId: 0
%% body:
如图所示，AO 是质量为 $m$ 的均匀细杆，可绕 O 轴在竖直平面内自由转动。细杆上的 P 点与放在水平桌面上的圆柱体接触，圆柱体靠在竖直的档板上而保持平衡。已知杆的倾角为 $\theta$，AP 长度是杆长的 $\frac{1}{4}$，各处的摩擦都不计，则档板对圆柱体的作用力等于 \tkanswer[2cm]{$\frac{1}{3}mg\sin2\theta$}。
\onepicture[align=right, w=4.5cm]{figs/25.png}

%% answer: $\frac{1}{3}mg\sin2\theta$ 或 $\frac{2}{3}mg\sin\theta\cos\theta$
\jdanswer{
$\frac{1}{3}mg\sin2\theta$ 或 $\frac{2}{3}mg\sin\theta\cos\theta$
}

%% memo:
\memoanswer{
由 AO 杆受力力矩平衡得 $mg\cdot\frac{l}{2}\cos\theta = N\cdot\frac{3l}{4}$，其中 $l$ 为杆长，解得 $N = \frac{2}{3}mg\cos\theta$；由题意知球受力平衡，只考虑水平方向，有 $F = \frac{2}{3}mg\sin\theta\cos\theta$。

\twopicture[w=4.5cm]{figs/25_an1.png}{figs/25_an2.png}
}

\item
%% number: 26
%% paperName: 1992年全国高考第 26 题 3 分
%% typeId: 4
%% type: 实验题
%% chapter: 电场
%% point: 描绘等势面
%% method: 无
%% score: 3
%% degree: 750
%% duplicateId: 0
%% body:
在用电流场模拟静电场描绘电场中等势线的实验中，所用的器材除了木板、白纸、复写纸、圆柱形电极、导线、电池、电键外，还必须有 \tkanswer[1.5cm]{导电纸}、\tkanswer[1.2cm]{探针} 和 \tkanswer[1.5cm]{电流表}。

%% answer: 导电纸，探针，电流表
\jdanswer{
导电纸；探针；电流表
}

%% memo:
\memoanswer{
本实验是用稳恒电流场模拟静电场，用电池、电流表、探针连成电路，在导电纸上找出等电势点，再描绘出等势线，还需用导电纸、探针和电流表。
}

\item
%% number: 27
%% paperName: 1992年全国高考第 27 题 3 分
%% typeId: 3
%% type: 填空题
%% chapter: 气体性质
%% point: 气体图象
%% method: 无
%% score: 3
%% degree: 450
%% duplicateId: 0
%% body:
图中直线 AB 为一定质量的理想气体等容过程的 $p-t$ 图线，原点 O 处的压强 $p=0$，温度 $t=0\UdC$。现先使该气体从状态 A 出发，经过一等温膨胀过程，体积变为原来体积的 2 倍，然后保持体积不变，缓慢加热气体，使之到达某一状态 F。此时其压强等于状态 B 的压强，试用作图方法，在所给的 $p-t$ 图上，画出 F 的位置。
\onepicture[align=right, w=3.5cm]{\includesvg{figs/27}}

%% answer: 如图所示
\jdanswer{
如图所示。只画出 F 点的位置，但未画出两条等容线相交于 $t$ 轴上一点者，不得分。

\onepicture[w=4.5cm]{\includesvg{figs/27_an}}
}

%% memo:
\memoanswer{
由状态 A 开始，经等温过程，且体积加倍，有

$p_{A}V_{A}=p_{C}V_{C}$，解得 $p_{C}=\frac{V_{A}}{V_{C}}p_{A}=\frac{p_{A}}{2}$，

由压强的关系确定 C 点，由于 $B\to A$ 是等容过程，反向延长 BA 与 $t$ 轴交于 D，D 应为热力学温度的 O 点，再连接 D、C 两点并延长到 F，使 BF 平行于 $t$ 轴，这个 F 点即为所求。
}

\item
%% number: 28
%% paperName: 1992年全国高考第 28 题 5 分
%% typeId: 6
%% type: 计算题
%% chapter: 光学
%% point: 透镜
%% method: 无
%% score: 5
%% degree: 750
%% duplicateId: 0
%% body:
一物体经焦距为 $24\Ucm$ 的凸透镜成一个放大率为 1.5 的实像。求物到透镜的距离。

%% answer: $u=40\Ucm$
\jdanswer{
$u=40\Ucm$
}

%% memo:
\memoanswer{
由题给数据根据透镜成像和放大率公式可得

$\frac{1}{u}+\frac{1}{v}=\frac{1}{f}=\frac{1}{24}$ ①

$m=\frac{v}{u}=1.5$ ②

解之得 $u=40\Ucm$ ③

评分标准：本题 5 分。得到①式给 2 分，得到②式给 2 分，求得最后结果，再给 1 分。
}

\item
%% number: 29
%% paperName: 1992年全国高考第 29 题 6 分
%% typeId: 6
%% type: 计算题
%% chapter: 电磁感应
%% point: 电磁感应受力分析
%% method: 无
%% score: 6
%% degree: 650
%% duplicateId: 0
%% body:
如图所示，导线框 abcd 固定在竖直平面内，bc 段的电阻为 $R$，其它电阻均可忽略。ef 是一电阻可忽略的水平放置的导体杆，杆长为 $l$，质量为 $m$，杆的两端分别与 ab 和 cd 保持良好接触，又能沿它们无摩擦地滑动。整个装置放在磁感应强度为 $B$ 的匀强磁场中，磁场方向与框面垂直。现用一恒力 $F$ 竖直向上拉 ef，当 ef 匀速上升时，其速度的大小为多少？
\onepicture[align=right, w=3cm]{figs/29.png}

%% answer: $v=\frac{R(F-mg)}{B^{2}l^{2}}$
\jdanswer{
$v=\frac{R(F-mg)}{B^{2}l^{2}}$
}

%% memo:
\memoanswer{
当杆 ef 向上运动时，杆中产生感应电动势。若杆向上运动的速度为 $v$，感应电动势为 $E = Blv$ ①

回路中的电流 $I = \frac{E}{R}$ ②

不论磁场的方向如何，安培力的方向总是向下。杆的平衡方程为 $F = BIl + mg$ ③

解以上 3 式得 $v=\frac{R(F-mg)}{B^{2}l^{2}}$ ④

评分标准：本题 6 分。求感应电动势占 1 分，求电流强度占 1 分，杆的平衡方程占 2 分，求得正确结果再给 2 分。
}

\item
%% number: 30
%% paperName: 1992年全国高考第 30 题 7 分
%% typeId: 6
%% type: 计算题
%% chapter: 气体性质
%% point: 玻意耳定律
%% method: 无
%% score: 7
%% degree: 450
%% duplicateId: 0
%% body:
如下左图所示，一个上下都与大气相通的直圆筒，内部横截面的面积 $S = 0.01\Umq$，中间用两个活塞 A 与 B 封住一定质量的理想气体，A、B 都可沿圆筒无摩擦地上、下滑动，但不漏气，A 的质量可不计、B 的质量为 $M$，并与一倔强系数 $k = 5\times10^{3}\UNm$ 的较长的弹簧相连。已知大气压强 $p_{0} = 1\times10^{5}\UPa$，平衡时，两活塞间的距离 $l_{0} = 0.6\Um$。现用力压 A，使之缓慢向下移动一定距离后，保持平衡。此时，用于压 A 的力 $F = 5\times10^{2}\UN$。求活塞 A 向下移的距离。（假定气体温度保持不变。）
\onepicture[align=right, w=2.8cm]{figs/30.png}

%% answer: $l=0.3\Um$
\jdanswer{
$l=0.3\Um$
}

%% memo:
\memoanswer{
活塞 A 受压向下移动的同时，活塞 B 也向下移动。已知达到平衡时，$F = 5\times10^{2}\UN$。设 A 向下移动的距离为 $l$，B 向下移动的距离为 $x$，由于气体温度不变，由玻意耳定律得

$p_{0}l_{0}S = \left(p_{0} + \frac{F}{S}\right)(l_{0} - l + x)S$ ①

当气体的压强为 $p_{0}$ 时，弹簧受 B 的作用而有一定的压缩量，当气体的压强变为 $p_{0} + \frac{F}{S}$ 时，弹簧增加的压缩量就是 B 向下移动的距离 $x$，由胡克定律得

$F = kx$ ②

由①、②两式消去 $x$，代入数字，得 $l = 0.3\Um$。

评分标准：本题 7 分。正确表示压缩后气体的压强、体积并列出①式，占 3 分；只写出玻意耳定律的普遍公式但未与此题所给各量联系起来的，不给这 3 分。通过文字说明或受力分析得到 B 移动的距离与 $F$ 的关系式②，占 3 分；只写出 $F = kx$ 而未说明 $x$ 代表什么的，不给这 3 分。求得最后结果再给 1 分。
}

\item
%% number: 31
%% paperName: 1992年全国高考第 31 题 8 分
%% typeId: 6
%% type: 计算题
%% chapter: 运动和力
%% point: 动量和能量
%% method: 无
%% score: 8
%% degree: 350
%% duplicateId: 0
%% body:
如图所示，一质量为 $M$、长为 $l$ 的长方形木板 B 放在光滑的水平地面上，在其右端放一质量为 $m$ 的小木块 A，$m < M$。现以地面为参照系，给 A 和 B 以大小相等、方向相反的初速度（如图），使 A 开始向左运动、B 开始向右运动，但最后 A 刚好没有滑离 B 板。以地面为参照系。

\begin{enumerate}
	\item
	若已知 A 和 B 的初速度大小为 $v_{0}$，求它们最后的速度的大小和方向。
	\item
	若初速度的大小未知，求小木块 A 向左运动到达的最远处（从地面上看）离出发点的距离。
\end{enumerate}
\onepicture[align=right, w=5cm]{\includesvg{figs/31}}

%% answer: 略
\jdanswer{
\begin{enumerate}
	\item
	$v = \frac{M - m}{M + m}v_{0}$，方向向右

	\item
	$l_{1}=\frac{M + m}{4M}l$
\end{enumerate}
}

%% memo:
\memoanswer{
（1）A 刚好没有滑离 B 板，表示当 A 滑到 B 板的最左端时，A、B 具有相同的速度。设此速度为 $v$，A 和 B 的初速度的大小为 $v_{0}$，则由动量守恒可得

$Mv_{0}-mv_{0}=(M+m)v$ ①

解得 $v=\frac{M-m}{M+m}v_{0}$，方向向右。

（2）A 在 B 板的右端时初速度向左，而到达 B 板左端时的末速度向右，可见 A 在运动过程中必经历向左作减速运动直到速度为零，再向右作加速运动直到速度为 $v$ 的两个阶段。设 $l_{1}$ 为 A 开始运动到速度变为零过程中向左运动的路程，$l_{2}$ 为 A 从速度为零增加到速度为 $v$ 的过程中向右运动的路程，$L$ 为 A 从开始运动到刚到达 B 的最左端的过程中 B 运动的路程，如图所示。设 A 与 B 之间的滑动摩擦力为 $f$，则由功能关系可知

对于 B：$fL=\frac{1}{2}Mv_{0}^{2}-\frac{1}{2}Mv^{2}$ ②

对于 A：$fl_{1}=\frac{1}{2}mv_{0}^{2}$ ③

$fl_{2}=\frac{1}{2}mv^{2}$ ④

由几何关系

$L+(l_{1}-l_{2})=l$ ⑤

由①、②、③、④、⑤式解得 $l_{1}=\frac{M+m}{4M}l$ ⑥

\onepicture[w=5cm]{figs/31_an.png}

评分标准：本题 8 分。

（1）2 分。末速度的大小和方向各占 1 分。

（2）6 分。其中关于 B 的运动关系式（例如②式）占 1 分；关于 A 的运动关系式（例如③、④两式）占 3 分，只要有错，就不给这 3 分；几何关系（例如⑤式）占 1 分；求出正确结果⑥，占 1 分。用其它方法求解，正确的，可参考上述评分标准进行评分。如考生若直接写出②、③、④、⑤的合并式

$fl=\frac{1}{2}(M+m)v_{0}^{2}-\frac{1}{2}(M+m)v^{2}$

则此式可给 2 分，再写出③式再给 3 分；最后结果正确再给 1 分。
}

\end{enumerate}

\end{document}
